


\section*{Begleitdokument für Lehrpersonen}
%\begin{itemize}
%    \item Einführung in den Semesterplan
%    \item Kontext (Voraussetzungen, Stufe, Empfehlung zur Einbettung und zur Fortsetzung)
%    \item Mögliche Schwierigkeiten der SuS und Lösungsvorschläge (für Gesamt-Programm)
%\end{itemize}

Dieses Dokument hat zum Ziel, einen Semesterplan für die Einführung in das Programmieren auf Gymnasialstufe\footnote{Schweizer Bildungssystem, Sekundarstufe II, Altersstufe ca. 15--18 Jahre, obligatorisches Fach Informatik, kurz: OInf} mittels Python zur Verfügung zu stellen. Das Semesterprogramm setzt keine besonderen Programmier- oder Informatikkenntnisse voraus. Das Semesterprogramm erstreckt sich über ca. 10 Lektionen und sollte den SuS ermöglichen, einfache mathematische und grafische Programme in Python zu schreiben.

Da relativ wenig Zeit für das OInf zur Verfügung steht, liegt der Fokus dieses Semesterprogramms auf grundlegenden Python-Kenntnissen. Als Fortsetzung und Vertiefung der neu gewonnen Kenntnisse empfiehlt sich eine Einführung in die Robotik im OInf, beziehungsweise eine Vertiefung der Programmierung im Ergänzungsfach, beispielsweise mit Themen wie Rekursion oder objektorientierter Programmierung (in Python).

Beim untenstehenden Themenüberblick wurde insbesondere darauf geachtet, die Reihenfolge möglichst sinnvoll zu gestalten, so dass die Kenntnisse schrittweise erweitert und aufeinander aufgebaut werden können. Zudem wurde darauf geachtet, stets ein möglichst ausgeglichenes Verhältnis zwischen \enquote{theoretisch-mathematischen} und \enquote{grafisch-interaktiven} Aufgaben zu bewahren (da letztere die SuS erfahrungsgemäss stärker motivieren).

\begin{overviewBox}[title=Überblick der Themen]
	\begin{enumerate}
		\item \textbf{Allgemeine Motivation und Einführung ins Programmieren}
		      \begin{itemize}
			      \item 1--2 Beispiele mit Bezug zur Lebenswelt und Vorwissen von SuS, einige Algorithmen und Beispiele zu logischem Denken
			      \item Vordefinierte Grafik-Befehle mit \lstinline|turtle| (nur ein Parameter): \lstinline|forward()|, \lstinline|right()|
			      \item Einfache Schleifen mit \lstinline|for i in range(...)|
		      \end{itemize}
		\item \textbf{Datentypen und Syntax}
		      \begin{itemize}
			      \item Variablen und Zuweisung vs. Mathematik
			      \item Arbeiten mit einfachen mathematischen Operatoren
			      \item Arbeiten mit Text
			      \item Debugging-Übungen (syntax vs. runtime vs. semantic errors), typische Fehler
		      \end{itemize}
		\item \textbf{Verzweigungen} (\lstinline|if...elif...else|) \& \lstinline|while|,  elementare logische Operatoren (or, and, xor)  und Flussdiagramme
		\item \textbf{Befehle / Funktionen} (\lstinline|def|)
		      \begin{itemize}
			      \item ohne Parameter, ohne \lstinline|return|
			      \item mit Parameter, ohne \lstinline|return|
			      \item mit Parameter, mit \lstinline|return|
		      \end{itemize}
		\item \textbf{Listen}
		      \begin{itemize}
			      \item  Koordinatengrafik, Zufallszahlen, \lstinline|random|-Library ($\rightarrow$ \lstinline|turtle|), \item Schleifen / Traversieren (\lstinline|for item in mylist|)
		      \end{itemize}
	\end{enumerate}
\end{overviewBox}


Mögliche Schwierigkeiten, die Auftreten können:
\begin{itemize}
	\item \textbf{Zeitbedarf}: Übungen zu neuen Teilen können erfahrungsgemäss viel Zeit beanspruchen, vornehmlich Übungen zu Syntax oder Funktionen. Dabei kann es hilfreich sein, die Übungen von Anfang an mit einer Zeiteinschätzung zu versehen und einzuplanen, welche Übungen bei Bedarf gestrichen, beziehungsweise auf eine andere Lektion verschoben, oder als Challenge-Aufgabe für Schnelle deklariert werden sollen.
	\item \textbf{Abstraktes Denken}: Gerade in Bereichen wie Verzweigungen und Flussdiagrammen, oder beim Traversieren von Listen, zeigt sich, dass das abstrakte Denken --- besonders in tieferen Altersstufen --- häufig noch nicht ausreichend entwickelt ist. Daher kann es, sofern diese Möglichkeit besteht, sinnvoll sein, das Semesterprogramm erst in einem späteren Semester auszuführen.
	\item \textbf{Technische Probleme}: \textit{Shortcuts}, Programme installieren, Programme abspeichern --- all dies ist für viele SuS entgegen der Erwartung vieler LP noch neu. Dazu kommen häufige Probleme beim Betreiben eigener Geräte (BYOD). Daher sollte genügend Zeit eingeplant werden, um auch auf solche Faktoren einzugehen und um \enquote{softes} Wissen zu vermitteln.
\end{itemize}

\clearpage
